\documentclass[aspectratio=169]{beamer}
\usetheme{Madrid}
\usecolortheme{beaver}
\usepackage[utf8]{inputenc}
\usepackage[T1]{fontenc}
\setbeamertemplate{navigation symbols}{}
\setbeamercolor{title}{fg=white,bg=cyan!60!black}
\setbeamercolor{frametitle}{fg=white,bg=cyan!50!black}
\setbeamercolor{block title}{bg=cyan!15,fg=black}
\setbeamercolor{block body}{bg=cyan!5,fg=black}

\title{G5 — Player Control}
\subtitle{Keyboard + mouse input schemes, aiming}
\author{ByteForge --- Grade 7 --- Game Maker}
\date{}

\begin{document}

\begin{frame}
    \titlepage
\end{frame}

\begin{frame}{Today's Plan}
    \begin{itemize}
        \item Move with Arrow Keys
        \item Aim with the Mouse
        \item Two-Player Split
        \item[$\bigstar$] BOSS --- Top-Down Shooter Controls
    \end{itemize}
\end{frame}

\begin{frame}{Keyboard input — four separate checks}
    \Large
    Checking all four arrow keys means four separate IF blocks — small changes every frame LOOK like smooth continuous movement.
\end{frame}


\begin{frame}{Mouse position — always available, live}
    \Large
    mouseX()/mouseY() always report the CURRENT position, live — nothing to set or update, the engine tracks it automatically.
\end{frame}


\begin{frame}{Two independent things, same frame loop}
    \Large
    Two players need two SEPARATE sets of variables — a single loop can control multiple independent things at once.
\end{frame}


\begin{frame}{BOSS — networking scale: PAN to LAN to WAN}
    \Large
    A keyboard+mouse on one PC is a PAN; school wifi is a LAN; the whole internet is a WAN — same idea, three different scales.
\end{frame}

\begin{frame}{Real-World Example}
    \Large
    A split-screen couch co-op game (two controllers, two characters on one screen) is exactly this chapter's "two independent input schemes, one frame loop" idea, running in a real shipped game.
\end{frame}

\begin{frame}{Watch}
    \Large
    Suggested video topic: \emph{Mouse and keyboard input in games explained for beginners}

    \vspace{1em}
    \normalsize
    Search this on YouTube together, or swap in a specific video you've
    already vetted --- do not project an unverified link cold in class.
\end{frame}

\begin{frame}{Let Them Discover}
    \Large
    Play any game with both keyboard movement AND mouse aiming (many shooters). Try to notice: does moving affect aiming, or are they genuinely independent? What would break if they shared one set of variables?
\end{frame}

\begin{frame}{Your Turn}
    \Large
    Open the student map, find \textbf{G5 — Player Control}, and work through it
    hands-on --- everything you just saw, they now do themselves.
\end{frame}

\end{document}
